\documentclass{article}
\usepackage[T1]{fontenc}
\usepackage[utf8]{inputenc}
\title{Incremental Heading Fixture}
\author{Fixture Author}
\date{\today}
\begin{document}

\maketitle

\section{Introduction}
This introductory section introduces the headings fixture. It contains the
distinctive term \texttt{headingfixture-zephyr}, which appears nowhere else in
the document, and is used by the tests to confirm that body text written under
a section heading is extracted faithfully. The paragraph is intentionally long
enough to wrap onto several lines so that extraction must reassemble wrapped
lines correctly.

\section{Core Findings}
A second top-level section follows the introduction and discusses the main
observations of the fixture. It provides a second standalone heading so that
tests can distinguish between multiple top-level sections rather than relying
on a single one.

\subsection{Methodology}
The methodology subsection describes how the fixture was produced. It repeats
the unique marker \texttt{headingfixture-zephyr} in a different paragraph so
that the content assertion remains robust whether or not the first occurrence
survives extraction intact.

\subsection{Results}
Results are summarized in plain prose here. Additional filler text about
measurement tolerances and instrument calibration ensures the subsection
contains enough content to be worth asserting on.

\section{Conclusion}
The conclusion wraps up the headings fixture and reiterates that the title and
section structure should be preserved in the extracted markdown so that heading
levels remain visible to downstream consumers.

\end{document}
